\subsection{How a physical defect bound becomes a count domain}
\label{ext:physical-count-domain}
This step explains the origin of the count coordinates and their joint
constraint. A small defect must be acquired for the same extremal law
before the following argument is used. Put $\sigma=\sqrt{10}$ and use
the physical defect $D=(3+\sigma)\rho-\gamma-3J$ from the conserved
identity, where $J$ is the gap sum.

Suppose $\supp P$ is bounded and has a central gap $(-A,B)$ with
$A,B>0$. Let $p=\Pp(X>0)$, $q=1-p$, $m=\E X_+=\E X_-$,
$V=\E(X\mid\ind_{\{X>0\}})$, and $W=X-V$. Write
\[
 \theta_+=\E[W^2;X>0],\quad \theta_-=\E[W^2;X<0],
 \quad\theta=\theta_++\theta_-,\quad
 \overline B=m/p,\quad\overline A=m/q.
\]
Then $m^2=(1-\theta)pq$. Set
$\Delta B=\overline B-B$ and $\Delta A=\overline A-A$; both
are nonnegative. The gap decomposition gives the joint inequality
\begin{align}
 D\ge{}&b(p,\theta)
 +6m\{(A+B)(\Delta A+\Delta B)+\Delta A\Delta B\}
 \notag\\
 &+3(\sigma+1)B\theta_++3(\sigma+3)A\theta_-,
 \label{ext:joint-physical-resource}\\
 b(p,\theta)={}&
 (6+2\sigma)(1-\theta)^{3/2}
       \frac{(p-p_*)^2}{\sqrt{pq}},\qquad p_*=(4-\sigma)/2.
 \notag
\end{align}
Here the binary charge, the inward displacement, and the conditional
variance all spend one common defect budget.

For clarity, the derivation needs only polynomial inequalities. Put
$a_\pm=\E X_\pm^2$, $b_\pm=\E X_\pm^3$, and
$T_+=b_+-mB^2$, $T_-=b_--mA^2$. Gap integration on each
side of the central gap gives
\[
 J-m(A+B)^2\le
 T_++T_--2B(a_+-mB)-2A(a_--mA).
\]
Substitute this bound in the definition of $D$. For a nonnegative
variable $Y$ of conditional mean $c$ one has exactly
\[
 \E Y^3=c^3+\E[(Y-c)^2(Y+2c)]
\]
under the conditional probability measure. Since $X\ge B$ on the
positive class and $-X\ge A$ on the negative class, this yields
$b_+\ge m\overline B^2+3B\theta_+$ and
$b_-\ge m\overline A^2+3A\theta_-$. Also
$a_+=m\overline B+\theta_+$ and
$a_-=m\overline A+\theta_-$. Substitution and collection of terms
give \eqref{ext:joint-physical-resource}; the constant part is
\[
 m\{(\sigma-1)\overline B^2+(\sigma+1)\overline A^2
          -6\overline A\overline B\}=b(p,\theta),
\]
because $(\sigma+1)p^2+(\sigma-1)q^2-6pq
=(6+2\sigma)(p-p_*)^2$.

Suppose a closed probability panel and a prior resource calculation
provide $p\in[l,u]\subset(0,1)$, $\theta\le\Theta<1$,
$A\ge A_l>0$, $B\ge B_l>0$, and $D\le d$.
Let $v_u=\max_{p\in[l,u]}p(1-p)$,
$k_+=3(\sigma+1)B_l$, $k_-=3(\sigma+3)A_l$, and let
$b_l\ge0$ be any proved lower bound for $b(p,\theta)$ on this
closed panel. These quantities may be replaced by rational bounds in
their indicated directions. With
\[
 \ell=\frac m{pq}=\sqrt{\frac{1-\theta}{pq}},\qquad
 X=\ell(\ind_{\{X>0\}}-p+U),
\]
put $d_Y=(\Delta A+\Delta B)/\ell$ and
$\lambda=n\E U^2=npq\theta/(1-\theta)$. Then
\begin{equation}\label{ext:count-triangle-derived}
 b_l+g d_Y+h\lambda\le D\le d,\qquad
 g=6(1-\Theta)(A_l+B_l),\quad
 h=\frac{\min(k_+,k_-)(1-\Theta)}{n v_u}.
\end{equation}
Indeed $m\ell=1-\theta$, so the displacement term in
\eqref{ext:joint-physical-resource} is at least $g d_Y$.
Its variance term is at least $\min(k_+,k_-)\theta\ge h\lambda$.
All coefficients are positive. Thus $d<b_l$ excludes the panel; if
$d=b_l$, the residual is deterministic; if $d>b_l$, the admissible
triangle is
\[
 \frac{d_Y}{d_0}+\frac\lambda\Lambda\le1,\qquad
 d_0=\frac{d-b_l}{g},\qquad\Lambda=\frac{d-b_l}{h}.
\]
The separate intercept bounds $d_Y\le d_0$ and
$\lambda\le\Lambda$ are weaker consequences of this one inequality.
They must not be used to invent an additional independent defect budget.

The construction is finite: each new moment, gap, or support bound must
follow from previously acquired simultaneous inequalities for this same
law. Intersecting those inequalities preserves the law. No assertion
about convergence of the refinement procedure is needed.
